\documentclass[UTF8,12pt,a4paper,fontset=fandol]{ctexart}

% 支持从项目根目录或本文件所在目录编译。
\makeatletter
\def\input@path{{../}{../../}}
\makeatother
\usepackage{researchnotes}

\title{行列式计算练习}
\author{}
\date{}

\begin{document}
\maketitle

\section*{作答说明}
本练习共 10 道行列式计算题，由基础计算逐步过渡到含参数及一般 $n$ 阶行列式。
请写出主要计算步骤，并将结果尽可能化简或写成因式的乘积。
除题目另有说明外，所有参数均为实数，$n$ 为整数。
本文件仅含题目。

\section*{计算题}
\begin{enumerate}[label=\arabic*.]

\item\label{q:detcalc-three}
计算
\[
\begin{vmatrix}
2&1&-1\\
1&3&2\\
3&-1&1
\end{vmatrix}.
\]

\item\label{q:detcalc-four}
计算
\[
\begin{vmatrix}
1&2&0&1\\
0&3&1&0\\
2&4&1&3\\
1&2&2&1
\end{vmatrix}.
\]

\item\label{q:detcalc-cyclic-three}
计算，并将结果分解因式：
\[
\begin{vmatrix}
a&b&c\\
c&a&b\\
b&c&a
\end{vmatrix}.
\]

\item\label{q:detcalc-parameter}
计算关于 $x$ 的行列式
\[
\begin{vmatrix}
x+1&x&x&x\\
x&x+2&x&x\\
x&x&x+3&x\\
x&x&x&x+4
\end{vmatrix}.
\]

\item\label{q:detcalc-powers}
计算
\[
\begin{vmatrix}
1&1&1&1\\
-2&-1&1&3\\
4&1&1&9\\
-8&-1&1&27
\end{vmatrix}.
\]

\item\label{q:detcalc-sparse}
计算
\[
\begin{vmatrix}
1&0&2&0\\
3&0&1&0\\
0&2&0&1\\
0&1&0&4
\end{vmatrix}.
\]

\item\label{q:detcalc-tridiagonal}
设 $n\geq2$，$A_n=(a_{ij})_{1\leq i,j\leq n}$，其中
\[
a_{ij}=\begin{cases}
3,&i=j,\\
1,&j=i+1,\\
2,&i=j+1,\\
0,&|i-j|\geq2.
\end{cases}
\]
计算 $D_n=\det A_n$，给出关于 $n$ 的通项公式。

\item\label{q:detcalc-min-powers}
设 $n\geq2$，计算
\[
D_n(a)=\det\bigl(a^{\min(i,j)}\bigr)_{1\leq i,j\leq n}.
\]
例如 $n=4$ 时，
\[
D_4(a)=\begin{vmatrix}
a&a&a&a\\
a&a^2&a^2&a^2\\
a&a^2&a^3&a^3\\
a&a^2&a^3&a^4
\end{vmatrix}.
\]
所得表达式应同时适用于 $a=0$ 和 $a=1$。

\item\label{q:detcalc-arrow}
设 $n\geq3$，$B_n(x)=(b_{ij})_{1\leq i,j\leq n}$，其中
\[
b_{ij}=\begin{cases}
x,&i=j,\\
1,&i=1,\ j\geq2,\text{ 或 }j=1,\ i\geq2,\\
0,&i\geq2,\ j\geq2,\ i\ne j.
\end{cases}
\]
计算 $\det B_n(x)$，并求其关于 $x$ 的全部实零点。

\item\label{q:detcalc-rational-vandermonde}
设 $n\geq2$，$x_1,\ldots,x_n$ 两两不同，且 $t\ne x_j$（$1\leq j\leq n$）。
定义 $n$ 阶矩阵 $C=(c_{ij})$：
\[
c_{ij}=\begin{cases}
x_j^{\,i-1},&1\leq i\leq n-1,\\
\displaystyle\frac{1}{t-x_j},&i=n.
\end{cases}
\]
其中零次幂表示常数多项式 $1$ 的取值。计算 $\det C$，将结果写成乘积之商。
例如 $n=3$ 时，所求行列式为
\[
\begin{vmatrix}
1&1&1\\
x_1&x_2&x_3\\
\dfrac{1}{t-x_1}&\dfrac{1}{t-x_2}&\dfrac{1}{t-x_3}
\end{vmatrix}.
\]

\end{enumerate}
\end{document}
